\item Un cilindro vertical con un pesado pistón contiene aire a 300 K. La presión inicial es $2.00 \times 10^5\text{ Pa}$ y el volumen inicial es $0.350\text{ m}^3$. Considere que la masa molar del aire es de $28.9\text{ g/mol}$ y suponga que $C_V = \frac{5}{2}R$. (a) Encuentre el calor específico del aire a volumen constante in unidades de $\text{J/kg}\cdot^\circ\text{C}$. (b) Calcule la masa del aire en el cilindro. (c) Suponga que el pistón se mantiene fijo. Determine la energía de entrada requerida para elevar la temperatura del aire a 700 K. (d) \textbf{¿Qué pasaría si?} Suponga de nuevo las condiciones del estado inicial y que el pesado pistón tiene libertad de movimiento. Encuentre la energía de entrada necesaria para elevar la temperatura a 700 K.
